\section{Pengenalan Library dan Framework Front-End}

Sub-CPMK 2.4 menuntut penggunaan library atau framework front-end untuk antarmuka dinamis \cite{mdn-learn}. \textbf{Bootstrap} menyediakan komponen siap pakai (navbar, card, form, modal) dan grid responsif; cukup sertakan CSS dan JS Bootstrap lalu gunakan class dan komponennya \cite{bootstrap}. \textbf{jQuery} menyederhanakan selektor DOM, event handling, dan AJAX dengan API yang ringkas \cite{jquery}. Untuk proyek yang lebih besar, framework seperti \textbf{Vue} atau \textbf{React} memungkinkan antarmuka reaktif dan komponen yang terorganisir \cite{mdn-learn}.

Dalam konteks mata kuliah, pengenalan Bootstrap atau jQuery sudah memadai: gunakan Bootstrap untuk layout dan komponen UI, atau jQuery untuk manipulasi DOM dan permintaan AJAX. Integrasi dengan back-end PHP (Bab X--XII) dapat dilakukan dengan \texttt{fetch} atau \texttt{\$.ajax} \cite{mdn-learn-js}.
